\ifthenelse{\equal{\lang}{ko}}{%
우리의 연구는 법률·정부 정보 접근, 다국어 NLP, 근거 기반 생성 및 책임 있는 AI 관련 기존 연구를 기반으로 하며,
기존 접근법의 한계를 드러내고 다국어·근거 기반 프레임워크의 참신성을 강조함으로써 기여를 위치시킨다.

\subsection{법률·정부 정보 접근}
공공 법률·정부 정보에 대한 접근을 개선하려는 초기 시도들은 키워드 기반 검색, 정적 포털, 정보 검색 시스템에 집중하였다 \cite{zhong2020legalai}.
법률 판단 예측(ECtHR) \cite{aletras2016echr}, 계약 이해 데이터셋(CUAD) \cite{hendrycks2021cuad} 등은 법률 텍스트의 자동 분석 가능성을 보여주었으나, 대부분 단일 언어이며 비전문가 대상 접근성보다는 전문가 과제에 초점을 둔다.
다국어 법률 코퍼스(MultiLegalPile) \cite{niklaus2023multilegal}는 저자원 언어로의 확장을 지원하지만, 시민 대면 정보 접근을 위한 통합 파이프라인은 제공하지 않는다.

\subsection{법률·정책 텍스트를 위한 다국어 NLP와 LLM}
Transformer 기반 인코더는 법률 NLP의 표준 도구가 되었다. LEGAL-BERT \cite{chalkidis2020legalbert}, LexGLUE 벤치마크 \cite{chalkidis2022lexglue}, LegalBench \cite{guha2023legalbench}는 법률 도메인 과제를 표준화하였다.
다국어 사전학습(XLM-R \cite{conneau2020xlmr}, mT5 \cite{xue2021mt5}) 및 대규모 다국어 번역(No Language Left Behind \cite{nllb2022})은 언어 간 전이를 통해 저자원 언어를 포괄하는 길을 열었고, 다국어 벤치마크(XTREME \cite{hu2020xtreme})는 이러한 능력을 체계적으로 평가한다.
GPT-4 \cite{openai2023gpt4}, LLaMA-2 \cite{touvron2023llama2}와 같은 LLM과 인간 피드백 기반 정렬 \cite{ouyang2022instructgpt}은 요약과 다국어 이해를 크게 향상시켰다.
그러나 이러한 능력을 비전문가 대상 공공 정보 접근에 통합한 연구는 여전히 드물다.

\subsection{근거 기반 생성과 책임 있는 AI}
사실성과 신뢰성을 강화하기 위한 연구가 활발하다. 검색 증강 생성(RAG) \cite{lewis2020rag}, 문장 임베딩 기반 의미 검색 \cite{reimers2019sbert, reimers2020multilingual}은 생성 결과를 외부 근거에 연결한다.
요약 충실성 \cite{maynez2020faithfulness}과 환각 탐지 \cite{ji2023hallucination}에 관한 연구는 LLM 출력 검증의 필요성을 강조한다.
동시에 Model Cards \cite{mitchell2021modelcards}, 언어 모델의 사회적 위험 분류 \cite{weidinger2021taxonomy, bender2021stochastic}, 그리고 EU AI 법안 \cite{EUAIAct2024}과 같은 규제는 투명성과 책임성의 중요성을 부각한다.
우리의 연구는 이러한 노력과 궤를 같이 하지만, 사실성 감사, 개념 드리프트 대응 \cite{gama2014drift}, 근거 연결을 배포 가능한 다국어 접근성 시스템에 직접 내장한다는 점에서 차별화된다.

\subsection{기여의 참신성(Novelty of Our Contribution)}
기존 연구와 달리, 우리는 법률·정부 정보 접근을 순수한 법률 전문가 과제로 취급하지 않는다.
대신, \textbf{변화 탐지–다국어 분류–근거 기반 요약·LLM 보고}를 결합한 \textbf{통합적이고 모듈형인 프레임워크}를 제안하며,
영어·한국어·프랑스어 데이터셋에서 검증하고 스페인어·아랍어·베트남어로 스트레스 테스트하였다.
이를 통해 우리는 비전문가와 다국어 공동체를 위한 \textbf{포용적 정보 접근}을 새로운 연구 방향으로 자리매김한다 \cite{devlin2019bert, bommasani2021foundation}.


}{%

Our work builds on prior research in legal and government information access, multilingual NLP, evidence-grounded generation, and responsible AI, situating our contribution by exposing the limitations of existing approaches and highlighting the novelty of our multilingual, evidence-grounded framework.


\subsection{Legal and Government Information Access}
Early efforts to improve access to public legal and government information focused on keyword-based search, static portals, and information-retrieval systems \cite{zhong2020legalai}.
Legal judgment prediction (ECtHR) \cite{aletras2016echr} and contract-understanding datasets (CUAD) \cite{hendrycks2021cuad} demonstrated the feasibility of automatically analyzing legal text, but most are monolingual and focus on expert tasks rather than accessibility for non-experts.
Multilingual legal corpora (MultiLegalPile) \cite{niklaus2023multilegal} support extension to low-resource languages, yet do not provide an integrated pipeline for citizen-facing information access.

\subsection{Multilingual NLP and LLMs for Legal and Policy Texts}
Transformer-based encoders have become standard tools in legal NLP. LEGAL-BERT \cite{chalkidis2020legalbert}, the LexGLUE benchmark \cite{chalkidis2022lexglue}, and LegalBench \cite{guha2023legalbench} standardized legal-domain tasks.
Multilingual pretraining (XLM-R \cite{conneau2020xlmr}, mT5 \cite{xue2021mt5}) and large-scale multilingual translation (No Language Left Behind \cite{nllb2022}) opened a path to covering low-resource languages via cross-lingual transfer, and multilingual benchmarks (XTREME \cite{hu2020xtreme}) systematically evaluate these capabilities.
LLMs such as GPT-4 \cite{openai2023gpt4} and LLaMA-2 \cite{touvron2023llama2}, together with human-feedback alignment \cite{ouyang2022instructgpt}, have substantially improved summarization and multilingual understanding.
However, work that integrates these capabilities into public information access for non-experts remains scarce.

\subsection{Evidence-Grounded Generation and Responsible AI}
A growing body of work strengthens factuality and reliability. Retrieval-augmented generation (RAG) \cite{lewis2020rag} and sentence-embedding-based semantic retrieval \cite{reimers2019sbert, reimers2020multilingual} link generated outputs to external evidence.
Studies on summarization faithfulness \cite{maynez2020faithfulness} and hallucination detection \cite{ji2023hallucination} emphasize the need to verify LLM outputs.
At the same time, Model Cards \cite{mitchell2021modelcards}, taxonomies of social risks from language models \cite{weidinger2021taxonomy, bender2021stochastic}, and regulations such as the EU AI Act \cite{EUAIAct2024} underscore the importance of transparency and accountability.
Our work aligns with these efforts but goes further by embedding factuality auditing, concept-drift adaptation \cite{gama2014drift}, and evidence linking directly into a deployable multilingual accessibility system.

\subsection{Novelty of Our Contribution}
In contrast to prior studies, we do not treat legal and government information access as a purely expert legal task.
Instead, we introduce a \textbf{unified, modular framework} that combines change detection, multilingual classification, and evidence-grounded summarization/LLM reporting,
validated on English, Korean, and French datasets and stress-tested on Spanish, Arabic, and Vietnamese.
By doing so, we position \textbf{inclusive information access} for non-experts and multilingual communities as a new research direction \cite{devlin2019bert, bommasani2021foundation}.

}
